\documentclass[hidelinks]{article}
\usepackage[a4paper, total={7in, 10in}]{geometry}
\usepackage[dvipsnames]{xcolor}
\usepackage{amsmath}
\usepackage{tikz}
\usepackage{tkz-euclide}
\usepackage[unicode]{hyperref}
\usepackage[all]{hypcap}
\usepackage{fancyhdr}
\usepackage{tocloft} % Add this package
\usepackage{xcolor}
\definecolor{darkgreen}{RGB}{0,100,0}
\usepackage{listings}
\lstset{
  language=R,
  basicstyle=\ttfamily\small,
  keywordstyle=\color{red},
  commentstyle=\color{darkgreen},
  stringstyle=\color{darkgreen},
  numbers=left,
  numberstyle=\tiny,
  stepnumber=1,
  numbersep=5pt,
  frame=single,
  breaklines=true,
  showstringspaces=false
}
\usepackage{amsmath}
\usepackage{amssymb}
\usepackage{pgfplots}
\pgfplotsset{compat=1.18}
\usepackage{amsmath}
\usepackage{float}
\usepackage{amsthm}
\usepackage[shortlabels]{enumitem}



% Remove dots from table of contents
\renewcommand{\cftdot}{}
\renewcommand{\cftsecdotsep}{\cftnodots}
\renewcommand{\cftsubsecdotsep}{\cftnodots}
\newcommand{\RNum}[1]{\uppercase\expandafter{\romannumeral #1\relax}}


\renewcommand{\contentsname}{Table of Contents}
\usetikzlibrary{angles,calc, decorations.pathreplacing}
\definecolor{carminered}{rgb}{1.0, 0.0, 0.22}
\definecolor{capri}{rgb}{0.0, 0.75, 1.0}
\definecolor{brightlavender}{rgb}{0.75, 0.58, 0.89}

\title{\textbf{STAT 545 HW 6}}
\author{Andres Efren Rocha Jayasinha}
\setcounter{secnumdepth}{0}
\date{Updated: \today} % suppresses the printed date

\begin{document}
\hypersetup{bookmarksnumbered=true,}
\pagecolor{white}
\color{black}
\maketitle
\begin{Large}
\tableofcontents
\end{Large}%
\pagebreak

\section{Assignment 6: Homogeneous Poisson Processes}

\subsection{Exercise 6.12}

\noindent Starting at noon, diners arrive at a restaurant according to a Poisson process at the rate of five customers per minute. The time each customer spends eating at the restaurant has an exponential distribution with mean $40$ minutes, independent of other customers and independent of arrival times. Find the distribution, as well as the mean and variance, of the number of diners in the restaurant at $2$ p.m.

\smallskip

\noindent \textbf{Solution}

\noindent Let $B_t$ denote the number of diners in the restaurant at time $t$ minutes after noon.
Arrivals occur according to a Poisson process with rate $\lambda=5$ per minute, and the
time spent eating by each diner is i.i.d.\ exponential with mean $40$ minutes. Thus the
eating time $Z$ has distribution $\mathrm{Exp}(\mu)$ with $\mu=1/40$ per minute.

\smallskip

\noindent We want the distribution of $B_{120}$.

\medskip

\noindent Let $N_t$ be the number of arrivals in $[0,t]$. Then $N_t\sim\mathrm{Poisson}(\lambda t)$.
Conditioning on $N_t$,
\begin{align*}
\mathbb{P}(B_t=k)
&=\sum_{n=k}^{\infty}\mathbb{P}(B_t=k\mid N_t=n)\,\mathbb{P}(N_t=n)\\
&=\sum_{n=k}^{\infty}\mathbb{P}(B_t=k\mid N_t=n)\,e^{-\lambda t}\frac{(\lambda t)^n}{n!}.
\end{align*}

\medskip
\noindent Given $N_t=n$, let $U_1,\dots,U_n$ be the (unordered) arrival times in $(0,t)$.
A standard property of the Poisson process states that conditional on $N_t=n$,
the variables $U_1,\dots,U_n$ are i.i.d.\ $\mathrm{Uniform}(0,t)$.

\smallskip

\noindent Let $Z_1,\dots,Z_n$ be the i.i.d.\ eating times, independent of $U_1,\dots,U_n$.
Diner $i$ is still in the restaurant at time $t$ iff $U_i+Z_i>t$. Hence $B_t$ equals the
number of indices $i\in\{1,\dots,n\}$ such that $U_i+Z_i>t$.

\smallskip

\noindent Define the indicator $I_i=\mathbf{1}\{U_i+Z_i>t\}$. Then, given $N_t=n$,
\[
B_t=\sum_{i=1}^n I_i.
\]
Moreover, conditional on $N_t=n$, the pairs $(U_i,Z_i)$ are i.i.d., so $I_1,\dots,I_n$ are
i.i.d.\ Bernoulli with success probability
\[
p=\mathbb{P}(U_1+Z_1>t).
\]
Therefore,
\[
\mathbb{P}(B_t=k\mid N_t=n)=\binom{n}{k}p^k(1-p)^{n-k}.
\]

\medskip
\noindent Since $U_1\sim \mathrm{Uniform}(0,t)$ and is independent of $Z_1$,
\begin{align*}
p
&=\mathbb{P}(U_1+Z_1>t)
=\frac{1}{t}\int_0^t \mathbb{P}(Z_1>t-x)\,dx.
\end{align*}
Here $Z_1\sim\mathrm{Exp}(\mu)$, so $\mathbb{P}(Z_1>y)=e^{-\mu y}$ for $y\ge 0$. Thus
\begin{align*}
p
&=\frac{1}{t}\int_0^t e^{-\mu(t-x)}\,dx
=\frac{1}{t}\int_0^t e^{-\mu u}\,du
=\frac{1-e^{-\mu t}}{\mu t}.
\end{align*}

\medskip
\noindent Substituting $\mathbb{P}(B_t=k\mid N_t=n)$ into the conditioning sum gives
\begin{align*}
\mathbb{P}(B_t=k)
&=\sum_{n=k}^{\infty}\binom{n}{k}p^k(1-p)^{n-k}e^{-\lambda t}\frac{(\lambda t)^n}{n!}\\
&=e^{-\lambda t}\frac{p^k(\lambda t)^k}{k!}
\sum_{n=k}^{\infty}\frac{\bigl((1-p)\lambda t\bigr)^{n-k}}{(n-k)!}\\
&=e^{-\lambda t}\frac{(p\lambda t)^k}{k!}\exp\!\bigl((1-p)\lambda t\bigr)\\
&=e^{-p\lambda t}\frac{(p\lambda t)^k}{k!},\qquad k=0,1,2,\dots
\end{align*}
Hence
\[
B_t\sim \mathrm{Poisson}(p\lambda t).
\]
Using $p=\dfrac{1-e^{-\mu t}}{\mu t}$, we get
\[
p\lambda t=\lambda\frac{1-e^{-\mu t}}{\mu}.
\]

\medskip
\noindent With $\lambda=5$, $\mu=1/40$, and $t=120$,
\[
p\lambda t=\frac{5}{1/40}\left(1-e^{-(1/40)\cdot 120}\right)=200(1-e^{-3}).
\]
Therefore,
\[
B_{120}\sim \mathrm{Poisson}\!\left(200(1-e^{-3})\right).
\]
For a Poisson$(\theta)$ random variable, $\mathbb{E}[B_{120}]=\mathrm{Var}(B_{120})=\theta$, so
\[
\mathbb{E}[B_{120}]=200(1-e^{-3}),
\qquad
\mathrm{Var}(B_{120})=200(1-e^{-3}).
\]

\subsection{Exercise 6.23}

\noindent Oak and maple trees are each located in the arboretum according to independent spatial Poisson processes with parameters $\lambda_O$ and $\lambda_M$, respectively.

\begin{enumerate}
    \item[(a)] In a region of $x$ square meters, find the probability that both species of trees are present.
    
    \item[(b)] In the arboretum, there is a circular pond of radius $100$ m. Find the probability that within $20$ m of the pond there are only oaks.
\end{enumerate}

{\color{darkgreen}

\noindent \textbf{Spatial Poisson Process}

\noindent A collection of random variables $(N_A)_{A \subseteq \mathbb{R}^d}$ is a spatial Poisson process with parameter $\lambda$ if

\begin{enumerate}
    \item For each bounded set $A \subseteq \mathbb{R}^d$, $N_A$ has a Poisson distribution with parameter $\lambda |A|$.
    
    \item Whenever $A$ and $B$ are disjoint sets, $N_A$ and $N_B$ are independent random variables.
\end{enumerate}
}

\noindent \textbf{Solution}

\noindent Part A:

\noindent Let $N_O$ be the number of oak trees and $N_M$ be the number of maple trees in a region of area $x$ (square meters).
Since the locations follow independent spatial Poisson processes with intensities $\lambda_O$ and $\lambda_M$,
\[
N_O \sim \mathrm{Poisson}(\lambda_O x), 
\qquad
N_M \sim \mathrm{Poisson}(\lambda_M x),
\]
and $N_O$ and $N_M$ are independent.

\noindent The event that \emph{both} species are present is
\[
\{N_O \ge 1,\; N_M \ge 1\}.
\]
Using complements and independence,

\[
P(O_x \ge 1,\, M_x \ge 1)
= P(O_x \ge 1)\, P(M_x \ge 1)
= (1 - P(O_x = 0))(1 - P(M_x = 0))
= (1 - e^{-\lambda_O x})(1 - e^{-\lambda_M x}).
\]

\noindent Part B:

\noindent Let $A$ be the region consisting of all points within $20$ m of the pond.
Since the pond has radius $100$ m, $A$ is an annulus with inner radius $100$ and outer radius $120$, so
\[
|A|=\pi(120^2-100^2).
\]
Let $O_A$ and $M_A$ denote the numbers of oaks and maples in $A$. By the spatial Poisson assumptions,
\[
O_A \sim \text{Poisson}(\lambda_O |A|), 
\qquad
M_A \sim \text{Poisson}(\lambda_M |A|),
\]
and $O_A$ and $M_A$ are independent.

\noindent The event ``within $20$ m of the pond there are only oaks'' means there are no maples in $A$, i.e. $\{M_A=0\}$.
Thus,

\[
P(O_R \ge 1,\; M_R = 0)
= P(O_R \ge 1)\,P(M_R=0)
= \bigl(1-e^{-\lambda_O 4400\pi}\bigr)e^{-\lambda_M 4400\pi}.
\]

\subsection{Exercise 6.35}

\noindent See the definitions for the spatial and nonhomogeneous Poisson processes.
Define a nonhomogeneous, spatial Poisson process in $\mathbb{R}^2$.
Consider such a process $(N_A)_{A \subseteq \mathbb{R}^2}$ with intensity function
\[
\lambda(x,y) = e^{-(x^2+y^2)}, 
\qquad -\infty < x,y < \infty.
\]

\noindent Let $C$ denote the unit circle, that is, the circle of radius $1$ centered at the origin.
Find $P(N_C = 0)$.

{\color{darkgreen}

\noindent \textbf{Nonhomogeneous Poisson Process}

\noindent A counting process $(N_t)_{t \ge 0}$ is a nonhomogeneous Poisson process with
intensity function $\lambda(t)$, if

\begin{enumerate}
    \item $N_0 = 0$.
    
    \item For all $t > 0$, $N_t$ has a Poisson distribution with mean
    \[
    \mathbb{E}(N_t) = \int_0^t \lambda(x)\,dx.
    \]
    
    \item For $0 \le q < r \le s < t$, 
    $N_r - N_q$ and $N_t - N_s$ are independent random variables.
\end{enumerate}

}

\noindent \textbf{Solution}

\noindent For a nonhomogeneous spatial Poisson process with intensity function
$\lambda(x,y)$, the number of points in a region $A$ has a Poisson
distribution with parameter
\[
\int\!\!\!\int_A \lambda(x,y)\,dx\,dy.
\]
Hence,
\[
P(N_A = 0)
= \exp\!\left(-\int\!\!\!\int_A \lambda(x,y)\,dx\,dy\right).
\]

\noindent In our problem, $C$ is the unit disk centered at the origin and
\[
\lambda(x,y)=e^{-(x^2+y^2)}.
\]
Therefore,
\[
P(N_C = 0)
=
\exp\!\left(-\int\!\!\!\int_C e^{-(x^2+y^2)}\,dx\,dy\right).
\]

\medskip

\noindent Since the region $C$ is a disk and the intensity depends only on
$x^2+y^2$, we switch to polar coordinates:
\[
x=r\cos\theta, \qquad y=r\sin\theta,
\]
with Jacobian determinant $r$.  On the unit disk,
\[
0 \le r \le 1, \qquad 0 \le \theta \le 2\pi.
\]

\noindent Thus,
\begin{align*}
\int\!\!\!\int_C e^{-(x^2+y^2)}\,dx\,dy
&=
\int_0^{2\pi} \int_0^1 e^{-r^2}\, r\,dr\,d\theta.
\end{align*}

\noindent First integrate with respect to $r$.  Let
\[
u = r^2, \qquad du = 2r\,dr,
\]
so that $r\,dr = \tfrac12 du$.  When $r=0$, $u=0$;
when $r=1$, $u=1$.  Hence,
\begin{align*}
\int_0^1 e^{-r^2} r\,dr
&=
\frac12 \int_0^1 e^{-u}\,du \\
&=
\frac12 \bigl(1 - e^{-1}\bigr).
\end{align*}

\noindent Now integrate with respect to $\theta$:
\begin{align*}
\int_0^{2\pi} \int_0^1 e^{-r^2} r\,dr\,d\theta
&=
\int_0^{2\pi} \frac12 (1-e^{-1})\, d\theta \\
&=
2\pi \cdot \frac12 (1-e^{-1}) \\
&=
\pi (1-e^{-1}).
\end{align*}

\medskip

\noindent Therefore,
\[
P(N_C = 0)
=
\exp\!\bigl(-\pi(1-e^{-1})\bigr).
\]

\subsection{Exercise 6.41}

\noindent Goals occur in a soccer game according to a Poisson process. 
The average total number of goals scored for a 90-minute match is 2.68. 
Assume that two teams are evenly matched. Use simulation to estimate the 
probability both teams will score the same number of goals. Compare with 
the theoretical result.

\noindent \textbf{Solution}

\noindent Since the total number of goals in 90 minutes follows a Poisson process 
with mean $2.68$, and the teams are evenly matched, each team’s number 
of goals is an independent Poisson random variable with mean $2.68/2 = 1.34$.

\noindent Let $X$ and $Y$ denote the number of goals scored by the two teams. Then
\[
X, Y \sim \text{Poisson}(1.34),
\quad \text{independently}.
\]

\noindent The probability that both teams score the same number of goals is
\[
P(X = Y)
=
\sum_{k=0}^{\infty}
P(X = k) P(Y = k)
=
\sum_{k=0}^{\infty}
\left(
\frac{e^{-1.34} 1.34^k}{k!}
\right)^2.
\]

\noindent Evaluating this sum gives
\[
P(X = Y) \approx 0.259.
\]

\noindent Thus, theoretically, the probability that both teams score the same 
number of goals is approximately $0.259$.

\noindent \textbf{Input:}

\begin{lstlisting}[caption={Poisson Soccer Simulation}]
trials <- 100000
mean(rpois(trials,1.34)==rpois(trials,1.34))
\end{lstlisting}

\noindent \textbf{Output:} We get .25812

\noindent \subsection{Exercise 6.42}
Simulate the restaurant results of Exercise 6.12.

\noindent \textbf{Solution}

\noindent \textbf{Input:}

\begin{lstlisting}[caption={South Carolina Swamp Forest}]
trials <- 10000
sim <- numeric(trials)
for (i in 1:trials) {
    cust <- rpois(1,5*60*2) # no. of customers in 2 hours
    arr <- sort(runif(cust,0,120))
    times <- arr + rexp(cust,1/40)
    sim[i] <- sum (times > 120) # no of customers at 2 pm
}

mean(sim)

var(sim)
\end{lstlisting}

\noindent \textbf{Output:} Mean of 190.1929 and Variance of 192.4839

\subsection{Exercise 6.43}

\noindent Simulate a spatial Poisson process with $\lambda = 10$ 
on the box of volume $8$ with vertices at the eight points 
$(\pm 1, \pm 1, \pm 1)$. Estimate the mean and variance of the number 
of points in the ball centered at the origin of radius $1$. 
Compare with the exact values.

\noindent \textbf{Solution}

\noindent For a spatial Poisson process with intensity $\lambda = 10$ per unit volume,
the number of points in any region $A$ follows a Poisson distribution with
mean
\[
\lambda \, |A|.
\]

\noindent The ball of radius $1$ has volume
\[
|A| = \frac{4}{3}\pi (1)^3 = \frac{4}{3}\pi.
\]

\noindent Hence the theoretical mean (and variance) of the number of points in the ball is
\[
10 \cdot \frac{4}{3}\pi = \frac{40\pi}{3} \approx 41.89.
\]

\noindent Thus, the exact distribution is
\[
N \sim \text{Poisson}\!\left(\frac{40\pi}{3}\right),
\]
so both the mean and variance equal $\frac{40\pi}{3} \approx 41.89$.

\noindent \textbf{Input:}

\begin{lstlisting}[caption={Exercise 6.43}]
trials <- 10000
sim <- numeric(trials)

for (k in 1:trials) {
  
  npoints <- rpois(1, lambda = 10 * 8)
  
  points <- matrix(runif(npoints * 3, -1, 1), ncol = 3)
  
  sim[k] <- sum(rowSums(points^2) < 1)
}

mean(sim)
var(sim)
\end{lstlisting}

\noindent \textbf{Output:} Mean of 41.9851 and Standard Deviation of 42.03528.

\end{document}